problem:  
Let \((X,d,m)\) be a compact \(\mathrm{RCD}(K,N)\) space with \(K>0\) and \(N<\infty\). By the Mondino–Naber and De Philippis–Gigli structure theorems, for \(m\)-almost every \(x\), there exists a unique Euclidean tangent cone \(T_xX \simeq \mathbb{R}^{k(x)}\), and the integer–valued dimension function \(k(x)\) is constant on each regular stratum.  
Denote by \(\mathcal{S}\subset X\) the singular set, and recall that \(\mathcal{S} = \bigcup_{j<N} \mathcal{S}^j\), where \(\mathcal{S}^j\) consists of those \(x\) whose tangents are \(\mathbb{R}^j\times C(Y)\) for some non‑Euclidean metric measure cone \(C(Y)\).  

**Problem:**  
Assume that for some \(j<N\), there exists a nonempty relatively open subset \(U\subset \overline{\mathcal{S}^j}\setminus\mathcal{S}^{j-1}\) such that the **volume density function**
\[
\vartheta(x) := \lim_{r\to 0}\frac{m(B_r(x))}{r^j}
\]
exists and is continuous on \(U\).  

Does this continuity of volume density along a fixed-dimensional singular stratum force the tangent cone to be unique at every point in \(U\)?  
Equivalently, can a continuous local measure profile coexist with *nonunique* tangent cones (with the same Euclidean factor but different cone cross-sections) in positively curved \(\mathrm{RCD}(K,N)\) spaces, or does continuity of \(\vartheta(x)\) impose rigidity of the cone cross-section \(Y\) across \(U\)?

Why is it a "good" problem:  
This question now sits directly at the current unknown frontier: the behavior of tangent cones and density functions along the **singular strata** of \(\mathrm{RCD}(K,N)\) spaces.  
While uniqueness and Euclidean regularity are understood for a.e. points, the geometry near singular strata remains poorly explored. This formulation isolates a sharp micro‑rigidity principle—does smooth variation of volume density enforce uniqueness of blow‑ups within a stratum?  

It is a "good" problem because:  
- It builds on rather than contradicts known structure theorems by limiting attention to the singular set, where open questions persist.  
- It forges a bridge between **metric‑measure analysis (density functions)** and **geometric stratification (cone uniqueness and cross‑section structure)**.  
- It is conceptually simple—continuity of volume versus rigidity of blow‑up—but potentially profound, touching upon the heart of synthetic curvature analysis, stability of tangent cones, and possible RCD analogues of geometric measure rigidity theorems.